% Generated draft: requires mathematical and source review.
\documentclass[11pt]{article}
\usepackage[margin=25mm]{geometry}
\usepackage{fontspec}
\setmainfont{Latin Modern Roman}
\usepackage{amsmath,amssymb,mathtools,unicode-math}
\setmathfont{Latin Modern Math}
\usepackage{enumitem,needspace,xurl,hyperref}
\setlength{\parindent}{0pt}
\setlength{\parskip}{5pt}
\setlength{\emergencystretch}{3em}
\newcommand{\N}{\mathbb{N}}
\newcommand{\Z}{\mathbb{Z}}
\newcommand{\Q}{\mathbb{Q}}
\newcommand{\R}{\mathbb{R}}
\newcommand{\C}{\mathbb{C}}
\newcommand{\F}{\mathbb{F}}
\newcommand{\A}{\mathbb{A}}
\newcommand{\PP}{\mathbb{P}}
\newcommand{\E}{\mathbb{E}}
\newcommand{\eps}{\varepsilon}
\newcommand{\dd}{\,\mathrm{d}}
\begin{document}
\section*{UnsolvedMath ID: 30001163}
\textbf{Source problem code:} catalogue entry
\subsection{Definitions and mathematical statement}
Let \(n\) denote the dimension parameter in the supplied formula, and let \(E_1\leq E_2\leq\cdots\) be the Dirichlet eigenvalues associated with a domain \(\Omega\). The supplied record does not specify the operator, the ambient space, the admissible class of domains, or hypotheses ensuring a discrete spectrum and well-defined real powers of the eigenvalues. The statement below is therefore conditional on the intended spectral setting, which requires clarification from the cited report.

For \(p>0\) and a positive integer \(J\), define
\[
M_p(J)=\left(\frac{n+2p}{nJ}\sum_{j=1}^{J}E_j^p\right)^{1/p}.
\]
The corrected formulation supplied in the record gives the inequality
\[
M_1(J)^2-M_2(J)^2\geq\frac{(E_{J+1}-E_J)^2}{4}.
\]
Its equality question is: does there exist an admissible domain \(\Omega\) and a finite positive integer \(J\) such that
\[
M_1(J)^2-M_2(J)^2=\frac{(E_{J+1}-E_J)^2}{4}?
\]
The admissibility conditions and whether \(n\) is fixed in advance remain unspecified in the supplied record.

The original extraction instead displayed \(M_{1/2}(J)-M_2(J)\) on the left. The supplied verification explicitly identifies this as a transcription error and reports that the official report uses \(M_1(J)^2-M_2(J)^2\). This draft follows that supplied correction without claiming independent source verification. The original extraction also mentions \(M_0(J)\), described as the corresponding geometric mean, without an explicit normalization; it is not needed for the corrected equality question.
\subsection{Short English statement}
The question asks whether the supplied universal inequality between normalized Dirichlet eigenvalue means has an exact equality case at a finite index. The draft follows the supplied correction to the squared means. The intended operator, domain hypotheses, and dimension quantifier still require clarification. The record labels the question open; that status is supplied dated curation, not an independently verified conclusion.
\subsection{Sources}
\begin{itemize}
\item Local UnsolvedMath registry, permanent ID \(30001163\), source problem code \texttt{catalogue entry}, title: Equality Cases for Universal Eigenvalue Mean Inequalities. The supplied record belongs to the Oberwolfach Reports Open Problems collection and gives the proposed year as 2009; no individual proposer is identified.
\item \emph{Low Eigenvalues of Laplace and Schrödinger Operators} (2009), Oberwolfach Reports, pp.\ 355--428, DOI 10.4171/owr/2009/06. Supplied source and literature reference: \url{https://doi.org/10.4171/owr/2009/06}. The supplied background describes these reports as documenting mathematical workshops and their open-problem discussions.
\item Harrell--Stubbe, universal bounds, as labeled in the supplied literature sources: \url{https://arxiv.org/abs/0808.1133}. No fuller bibliographic description is supplied.
\item Supplied statement verification: the record labels the statement \texttt{corrected\_verified} and reports that comparison with the official report corrected the left side from \(M_{1/2}(J)-M_2(J)\) to \(M_1(J)^2-M_2(J)^2\). This verification is reported from the record, not independently performed here.
\item Literature assessment checked at \texttt{2026-08-21}: the supplied curation labels the problem open, describes the trace-identity inequality as known and asymptotically sharp, and reports that no domain and index giving exact finite-index saturation was found. This is supplied dated curation derived from the OWR record, not independently verified; the status and cited mathematical claims require independent checking.
\end{itemize}
\end{document}
